\documentclass[11pt]{article}
\usepackage[margin=1in]{geometry}
\usepackage[T1]{fontenc}
\usepackage{lmodern}
\usepackage{amsmath,amssymb}
\setlength{\parindent}{0pt}
\setlength{\parskip}{0.8em}
\title{Linear Algebra: Systems, Row Reduction, Rank, and Determinants}
\author{Intermediate Course Notes}
\date{}
\begin{document}
\maketitle

\section*{Systems of Linear Equations}

A system of $m$ linear equations in $n$ unknowns $x_1, x_2, \ldots, x_n$ has the form:
\[
\begin{aligned}
a_{11} x_1 + a_{12} x_2 + \cdots + a_{1n} x_n &= b_1, \\
a_{21} x_1 + a_{22} x_2 + \cdots + a_{2n} x_n &= b_2, \\
&\vdots \\
a_{m1} x_1 + a_{m2} x_2 + \cdots + a_{mn} x_n &= b_m.
\end{aligned}
\]
In matrix form, this is $Ax = b$, where $A$ is the $m \times n$ coefficient matrix, $x$ is the $n \times 1$ column vector of unknowns, and $b$ is the $m \times 1$ column vector of constants.

A solution is a vector $x$ that satisfies all $m$ equations simultaneously. There are exactly three possibilities:
\begin{enumerate}
  \item Exactly one solution (the system is consistent and the solution is unique).
  \item No solution (the system is inconsistent: the equations contradict each other).
  \item Infinitely many solutions (the system is consistent but underdetermined: some variables are free).
\end{enumerate}

Geometrically, each equation defines a hyperplane in $\mathbb{R}^n$. Solutions are the points where all hyperplanes intersect. In $\mathbb{R}^2$, two equations typically define two lines, and a unique solution is their intersection point. If the lines are parallel, there is no intersection. If they are the same line, every point on it is a solution.

The augmented matrix $[A \mid b]$ encodes the entire system in a single object:
\[
[A \mid b] = \begin{bmatrix} a_{11} & a_{12} & \cdots & a_{1n} & b_1 \\ a_{21} & a_{22} & \cdots & a_{2n} & b_2 \\ \vdots & \vdots & & \vdots & \vdots \\ a_{m1} & a_{m2} & \cdots & a_{mn} & b_m \end{bmatrix}.
\]

\section*{Gaussian Elimination and Row Reduction}

Gaussian elimination is the systematic method for solving $Ax = b$ by transforming the augmented matrix into row echelon form. The three elementary row operations are:
\begin{itemize}
  \item \textbf{Row swap}: interchange rows $i$ and $j$.
  \item \textbf{Row scaling}: multiply row $i$ by a nonzero scalar $c$.
  \item \textbf{Row addition}: add $c$ times row $j$ to row $i$.
\end{itemize}
None of these operations change the solution set of the system.

A matrix is in row echelon form when:
\begin{itemize}
  \item All zero rows are at the bottom.
  \item Each nonzero row's leading entry (pivot) is strictly to the right of the pivot in the row above.
\end{itemize}

A matrix is in reduced row echelon form (RREF) when additionally:
\begin{itemize}
  \item Each pivot equals 1.
  \item Each pivot is the only nonzero entry in its column.
\end{itemize}

Example: solve the system
\[
\begin{cases} x + 2y = 5 \\ 2x - y = 0. \end{cases}
\]
Augmented matrix:
\[
\begin{bmatrix} 1 & 2 & 5 \\ 2 & -1 & 0 \end{bmatrix}
\xrightarrow{R_2 \leftarrow R_2 - 2R_1}
\begin{bmatrix} 1 & 2 & 5 \\ 0 & -5 & -10 \end{bmatrix}
\xrightarrow{R_2 \leftarrow -\frac{1}{5}R_2}
\begin{bmatrix} 1 & 2 & 5 \\ 0 & 1 & 2 \end{bmatrix}
\xrightarrow{R_1 \leftarrow R_1 - 2R_2}
\begin{bmatrix} 1 & 0 & 1 \\ 0 & 1 & 2 \end{bmatrix}.
\]
Reading off RREF: $x = 1$, $y = 2$.

For a system with free variables, the RREF will have columns without pivots. Those columns correspond to free variables that can take any value, generating infinitely many solutions.

\section*{Identifying Consistency from Row Echelon Form}

After row reducing the augmented matrix, check whether the system is consistent:
\begin{itemize}
  \item If there is a pivot in the augmented column (a row of the form $[0 \; 0 \; \cdots \; 0 \mid c]$ with $c \neq 0$), the system is inconsistent: no solution.
  \item If no pivot appears in the augmented column, the system is consistent.
    \begin{itemize}
      \item If every variable column has a pivot, the solution is unique.
      \item If some variable columns lack pivots (free variables), there are infinitely many solutions.
    \end{itemize}
\end{itemize}

This analysis makes row reduction not just a computational tool but a structural diagnostic: it tells you everything about whether and how a system can be solved.

\section*{Rank of a Matrix}

The rank of a matrix $A$ is the number of pivot positions in any row echelon form of $A$. It equals the dimension of the column space (image) of $A$, and also the dimension of the row space.

Formally, the column space $\text{Col}(A)$ is the span of the columns of $A$. The rank counts how many columns are genuinely independent (i.e., not expressible as a combination of earlier columns).

Key facts about rank:
\begin{itemize}
  \item $\text{rank}(A) \leq \min(m, n)$ for an $m \times n$ matrix.
  \item $A$ has full column rank if $\text{rank}(A) = n$: all columns are independent. The equation $Ax = b$ has at most one solution for each $b$.
  \item $A$ has full row rank if $\text{rank}(A) = m$: every $b$ is achievable. The equation $Ax = b$ always has at least one solution.
  \item A square $n \times n$ matrix has full rank $n$ if and only if it is invertible.
\end{itemize}

The null space (kernel) of $A$ is $\ker(A) = \{x : Ax = \mathbf{0}\}$. The rank-nullity theorem states:
\[
\text{rank}(A) + \text{nullity}(A) = n,
\]
where $n$ is the number of columns. This says the dimensions of the image and kernel are complementary: each pivot variable corresponds to a dimension of the image, and each free variable corresponds to a dimension of the kernel.

Example: the matrix
\[
A = \begin{bmatrix} 1 & 2 & 3 \\ 2 & 4 & 6 \end{bmatrix}
\]
row-reduces to
\[
\begin{bmatrix} 1 & 2 & 3 \\ 0 & 0 & 0 \end{bmatrix},
\]
which has rank 1. Nullity is $3 - 1 = 2$: there is a two-dimensional family of vectors mapped to zero.

\section*{Linear Independence via Row Reduction}

Row reduction provides the definitive test for linear independence. Given vectors $v_1, \ldots, v_k \in \mathbb{R}^m$, form the matrix $A$ with these as columns. Row-reduce $A$:
\begin{itemize}
  \item If every column of $A$ has a pivot, the vectors are linearly independent.
  \item If any column lacks a pivot, that column's vector is a linear combination of the previous ones, and the set is dependent.
\end{itemize}

The pivot columns of $A$ form a basis for the column space. Non-pivot columns are in the span of the pivot columns to their left.

This also connects to systems: $c_1 v_1 + \cdots + c_k v_k = \mathbf{0}$ is a homogeneous system $Ac = \mathbf{0}$. The vectors are independent iff this system has only the trivial solution, which happens iff every column of $A$ is a pivot column.

\section*{Determinants: Definition and Computation}

The determinant $\det(A)$ is a scalar defined for square matrices. For $1 \times 1$: $\det([a]) = a$. For $2 \times 2$:
\[
\det\begin{bmatrix} a & b \\ c & d \end{bmatrix} = ad - bc.
\]

For larger matrices, the determinant is computed by cofactor expansion along any row or column. Expanding along row $i$:
\[
\det(A) = \sum_{j=1}^n (-1)^{i+j} A_{ij} \det(M_{ij}),
\]
where $M_{ij}$ is the $(n-1) \times (n-1)$ minor obtained by deleting row $i$ and column $j$, and $(-1)^{i+j}$ is the sign factor (the ``checkerboard'' of $+$ and $-$ signs).

For a $3 \times 3$ matrix, expanding along the first row:
\[
\det\begin{bmatrix} a & b & c \\ d & e & f \\ g & h & i \end{bmatrix} = a(ei - fh) - b(di - fg) + c(dh - eg).
\]

Properties of the determinant under row operations:
\begin{itemize}
  \item Swapping two rows multiplies the determinant by $-1$.
  \item Multiplying a row by $c$ multiplies the determinant by $c$.
  \item Adding a multiple of one row to another does not change the determinant.
\end{itemize}

Using these properties, one can row-reduce a matrix to triangular form and compute the determinant as the product of the diagonal entries (adjusted for any row swaps and row scalings performed along the way).

The determinant of a triangular matrix (upper or lower) is the product of its diagonal entries.

\section*{Geometric Meaning of the Determinant}

The absolute value $|\det(A)|$ measures how much the linear transformation $x \mapsto Ax$ scales volumes. More precisely:
\begin{itemize}
  \item In $\mathbb{R}^2$: $|\det(A)|$ is the area of the parallelogram spanned by the columns of $A$.
  \item In $\mathbb{R}^3$: $|\det(A)|$ is the volume of the parallelepiped spanned by the columns.
  \item In general: $|\det(A)|$ is the $n$-dimensional volume scaling factor.
\end{itemize}

The sign of $\det(A)$ indicates orientation:
\begin{itemize}
  \item $\det(A) > 0$: the transformation preserves orientation.
  \item $\det(A) < 0$: the transformation reverses orientation (like a reflection).
  \item $\det(A) = 0$: the transformation is singular; it collapses the space to a lower dimension and destroys volume.
\end{itemize}

The product rule for determinants: $\det(AB) = \det(A) \det(B)$. This says the volume scaling of a composed transformation is the product of the individual scalings.

\section*{Determinants and Invertibility}

A square matrix $A$ is invertible if and only if $\det(A) \neq 0$.

When $\det(A) = 0$, the matrix is singular: its columns are linearly dependent, its rank is less than $n$, and the transformation $x \mapsto Ax$ collapses the input space to a lower-dimensional image. There is no way to recover the original input from the output.

When $\det(A) \neq 0$, the transformation is bijective: every output corresponds to exactly one input, and the inverse transformation exists.

Cramer's rule gives an explicit formula for the solution of $Ax = b$ when $A$ is invertible:
\[
x_j = \frac{\det(A_j)}{\det(A)},
\]
where $A_j$ is the matrix obtained by replacing column $j$ of $A$ with $b$. Cramer's rule is elegant but computationally expensive for large systems; Gaussian elimination is preferred in practice.

\section*{The Matrix Inverse via Row Reduction}

To compute $A^{-1}$ for an invertible $n \times n$ matrix $A$: augment $A$ with the identity matrix $I$ and row-reduce the entire augmented matrix $[A \mid I]$ to RREF. When the left block becomes $I$, the right block is $A^{-1}$:
\[
[A \mid I] \xrightarrow{\text{RREF}} [I \mid A^{-1}].
\]

Example: let $A = \begin{bmatrix} 2 & 1 \\ 5 & 3 \end{bmatrix}$. Then $\det(A) = 6 - 5 = 1 \neq 0$, so $A$ is invertible.

\[
\begin{bmatrix} 2 & 1 & 1 & 0 \\ 5 & 3 & 0 & 1 \end{bmatrix}
\to
\begin{bmatrix} 1 & 0 & 3 & -1 \\ 0 & 1 & -5 & 2 \end{bmatrix}.
\]
So $A^{-1} = \begin{bmatrix} 3 & -1 \\ -5 & 2 \end{bmatrix}$.

Check: $AA^{-1} = \begin{bmatrix} 2 & 1 \\ 5 & 3 \end{bmatrix}\begin{bmatrix} 3 & -1 \\ -5 & 2 \end{bmatrix} = \begin{bmatrix} 1 & 0 \\ 0 & 1 \end{bmatrix}$. Confirmed.

Properties of the inverse: $(A^{-1})^{-1} = A$; $(AB)^{-1} = B^{-1}A^{-1}$; $(A^T)^{-1} = (A^{-1})^T$; $\det(A^{-1}) = 1/\det(A)$.

\section*{Vector Spaces and Subspaces}

A vector space $V$ over $\mathbb{R}$ is a set equipped with addition and scalar multiplication satisfying eight axioms: commutativity and associativity of addition, the existence of a zero vector, the existence of additive inverses, and four distributive/associative laws connecting addition and scalar multiplication.

Common examples: $\mathbb{R}^n$ (tuples of $n$ real numbers), $\mathcal{P}_n$ (polynomials of degree at most $n$), $\mathcal{M}_{m \times n}$ (matrices of size $m \times n$), and the zero space $\{0\}$.

A subspace of $V$ is a subset $W \subseteq V$ that is itself a vector space. The three subspace conditions: $\mathbf{0} \in W$; if $u, v \in W$ then $u + v \in W$; if $u \in W$ and $c \in \mathbb{R}$ then $cu \in W$.

Important subspaces associated with a matrix $A$ of size $m \times n$:
\begin{itemize}
  \item Column space $\text{Col}(A) \subseteq \mathbb{R}^m$: span of the columns of $A$.
  \item Null space $\text{Null}(A) \subseteq \mathbb{R}^n$: $\{x : Ax = \mathbf{0}\}$.
  \item Row space $\text{Row}(A) \subseteq \mathbb{R}^n$: span of the rows of $A$.
\end{itemize}

The dimension of $\text{Col}(A)$ and the dimension of $\text{Row}(A)$ are both equal to $\text{rank}(A)$. Row reduction preserves the row space but may change the column space; use the original pivot columns to form a basis for the column space.

\section*{Eigenvalues and the Characteristic Polynomial}

For an $n \times n$ matrix $A$, the characteristic polynomial is
\[
p(\lambda) = \det(A - \lambda I).
\]
This is a degree-$n$ polynomial in $\lambda$. The eigenvalues of $A$ are the roots of $p(\lambda) = 0$.

Once an eigenvalue $\lambda$ is known, the corresponding eigenvectors are found by solving the homogeneous system $(A - \lambda I)x = \mathbf{0}$. The set of solutions is the eigenspace $E_\lambda = \ker(A - \lambda I)$, which is a subspace.

The algebraic multiplicity of an eigenvalue $\lambda$ is its multiplicity as a root of the characteristic polynomial. The geometric multiplicity is $\dim(E_\lambda) = \text{nullity}(A - \lambda I)$.

The geometric multiplicity is always at most the algebraic multiplicity. When they are equal for every eigenvalue, the matrix is diagonalizable: $A = PDP^{-1}$ where $P$ contains eigenvectors as columns and $D$ is diagonal with eigenvalues on the diagonal.

If all $n$ eigenvalues are distinct, the matrix is automatically diagonalizable.

\section*{Row Reduction and Eigenvalues}

Row reduction should never be used to find eigenvalues: row operations change the determinant (and thus the characteristic polynomial) and do not preserve eigenvalues. Always compute eigenvalues from $\det(A - \lambda I) = 0$ using the original matrix.

Row reduction is appropriate for: solving $Ax = b$, finding $A^{-1}$, computing rank and null space, and testing linear independence. Eigenvalues require the characteristic polynomial, not row echelon form.

The trace of $A$ (sum of diagonal entries) equals the sum of eigenvalues. The determinant of $A$ equals the product of eigenvalues. These are fast checks when the eigenvalues are already known.

\section*{Summary}

Row reduction is the central computational tool of this course: it solves systems, computes rank and null space, finds inverses, and tests independence. Rank measures the genuinely independent content of a matrix, and the rank-nullity theorem constrains how much information can be preserved or lost.

Determinants measure volume scaling and detect singularity. Together, rank, determinants, and inverses answer the most basic question of linear algebra: when does a linear system have solutions, and when do those solutions exist uniquely?

\end{document}
